\section{An equal middle and final run in the tail}\label{tail:diagonal-sec}
\begin{theorem}\label{tail:diagonal}
For every $k\ge1$ and $0\le i<k$, the triple $(i,k,k)$ forces
periodicity without a length condition.
\end{theorem}
\begin{proof}
Orient a solution with unequal corresponding lengths so $d=x-u>0$.
Put $p=(k-i)d/3>0$. This is an integer since $D=(i+2k)d/3=kd-p$
is an integer. With $\ell=k$, the hypotheses of the factorization
\eqref{eq:tail} hold: $id\le D$ and $i+k<2k$.
In its notation $|G|=p$, $|A|=2p$, $|T'|=D$, and $|H|=D-p$.
The words $H$ and $T'$ are prefixes of the very same $X^k$:
for $H$ use $X^k=HAU^k$, and for $T'$ use $X^kG=T'AU^k$
and $D<kx$. Write $T'=HK$, where $|K|=p$.
Then
\[
 ACY=YT'=YHK=GCYK .
\]
Comparing the initial $p$ letters gives $A=GJ$ with $|J|=p$.
Cancelling $G$ gives $JCY=CYK$.
A second use of the equality of the two run exponents gives
\[
 HAU^kG=X^kG=HKAU^k,\qquad KAU^k=AU^kG .
\]
In particular $A$ is the prefix of $KA$ of length $2p$.
Its first $p$ letters are $K$ and its next $p$ letters repeat them.
Hence $A=KK$, and $A=GJ$ gives $G=J=K$.
Cancelling $A=K^2$ shows $KU^k=U^kK$; also $KCY=CYK$.
As $K$ is nonempty, the commutation theorem puts $U$ and
$V=CYA$ in the powers of the primitive root of $K$, including
the possibility of empty images. Thus the right morphism is periodic.
Lemma~\ref{perred} finishes: family A has equal internal exponents,
family B has zero middle exponent here, and the single-$a$ sporadics
have $S=1$, whereas $S=i+2k\ge2$.
\end{proof}

\begin{lemma}\label{tail:112}
For $(1,1,2)$ write $d=3e$. A counterexample must have $u<e$.
\end{lemma}
\begin{proof}
Suppose otherwise. We can assume $XU\ne UX$ by Lemma~\ref{commouter},
so $e\le u<3e$ by Corollary~\ref{outerbound}. In \eqref{eq:tail}
the lengths of $A,G,C,T',H$ are respectively $e,2e,3e,4e,5e$.
Since $4e\le x$, both $T'$ and $H$ are prefixes of $X^\omega$,
and $H=T'K$ with $|K|=e$. Now $GCY=YH=ACYK$ gives $G=AJ$
and $JCY=CYK$, $|J|=e$. The latter equation makes $CY$ a prefix
of $J^\omega$, by iteration; thus $C=J^3$.
The word $X=UC$ ends with $U$, so $u<3e$ makes $U$ a suffix of $C$.
Consequently $X$ is a factor of $J^{\mathbb Z}$ ending at phase zero.
Its prefix $T'$ has length $4e$, hence is a fourth power of a rotation
of $J$. Its conjugate $AC=AJ^3$ is likewise a fourth power of a
length-$e$ word. Comparing the four blocks shows $A=J$ and $G=J^2$.
Finally $XG=T'AU$ reads $UJ^5=T'JU$. On the left the whole word
has period the primitive root length $r$ of $J$ and starts at phase
$-u$. On the right a full $J$ starts at offset $4e$. Primitivity
forces $-u+4e\equiv0\pmod r$, so $r\mid u$. The suffix $U$ of
$J^3$ therefore commutes with $J$ and $X$, a contradiction.
\end{proof}

\begin{lemma}\label{tail:013}
For $(0,1,3)$ put $d=3e$. If $y+\min(x,4e)\ge8e$, all images
commute.
\end{lemma}
\begin{proof}
Assume $XU\ne UX$, so $0<u<3e$ by the outer bound and the
periodic-morphism reduction. Write $A,G$ for the words of lengths
$4e,5e$ in \eqref{eq:tail}. Here $C$ is empty,
$AY=YT'$, $GY=YH$, $XG=T'AU$, and $X^3=HAU^3$.
Thus $Y$ is a common prefix of $A^\omega,G^\omega$. After this
prefix those two infinite words continue with $T'$ and $H$
respectively. These two continuations agree for at least
$\min(x,4e)$ letters: $T'$ starts the word $XG$ and $H$ starts
$X^3$. Fine--Wilf on the common prefix gives period $e$, so
$A=P^4$, $G=P^5$ where $|P|=e$. Write $r$ for the primitive
root length of $P$. The equation $XG=T'AU$ ends in $U$, so
$U$ is a suffix of $P^5$. In this same equation the two
$P$-periodic intervals $[x,x+5e)$ and $[4e,8e)$ start at phase zero.
Their overlap has length at least $2e$, since $3e<x<6e$ and
$x+5e>8e$. Equality of a full primitive rotation gives $r\mid x-4e$,
hence $r\mid u$. It follows that $AU^3$ is a power of the primitive
root of $P$. As a suffix of $X^3$ it has periods $r,x$ and length
$4e+3u\ge x+r$. Fine--Wilf makes this suffix periodic with period
$\gcd(r,x)$, which divides $x$ and $u$. Its terminal length-$x$
and length-$u$ words are $X$ and $U$ at the same ending phase.
They commute, giving the desired contradiction.
\end{proof}

\begin{corollary}
In the remaining range $u<e$ of Lemma~\ref{tail:112}, a
counterexample also requires $y<2e-u$.
\end{corollary}
\begin{proof}
Align the two bi-infinite repetitions of $AC$ and $GC$ at the starts
of these written blocks. Immediately to the left they agree on $C$
of length $3e$. To the right they agree on $Y$, after which they
start with $T'$ and $H$. These agree for the first $x$ positions,
since $x<4e$ and both start with $X$. The common interval thus has
length at least $3e+y+x$. If $y\ge2e-u$ this is at least $8e$,
the Fine--Wilf threshold for $4e,5e$. Consequently the two entire
bi-infinite repetitions coincide with period $e$: extend the common
$e$-periodic interval using the periods $4e$ and $5e$, which are
multiples of $e$ and each at most its length. At the common origin
this yields $A=J,G=J^2,C=J^3$ for one length-$e$ word $J$.
The final phase argument in Lemma~\ref{tail:112} now applies verbatim:
$U$ is a suffix of $J^3$ and $UJ^5=T'JU$, so the primitive root
length of $J$ divides $u$.
\end{proof}

\begin{theorem}[An exponent criterion on either side of the band]\label{tail:windows}
Assume $i+k<2\ell$, $k\ge1$, and define integers
\[
 c=i+\ell-2,\quad m=i+\ell-2k\ne0,\quad
 a=\min(3i,i+k+\ell),\quad
 b=3k+\min(0,2m),\quad s=3k+\max(m,2m).
\]
If $c>0$, $|m|\le c+1$, and $c(a+b-s)\ge3(c+1)$,
the tuple forces periodicity.
\end{theorem}
\begin{proof}
In a counterexample all words are nonempty by Lemma~\ref{perred}
and $r>(c+1)u$, $u<d/c$ by Corollary~\ref{outerbound}.
We construct two bi-infinite words, of periods $D$ and
$h=(i+\ell)d-D$, agreeing on three consecutive pieces
$L,Y,F$. Here $L$ is a suffix of $X^i$ of length $ad/3$;
$F$ is a prefix of $X^\omega$ of length $bd/3$.
The longer period is $Q=sd/3$, the shorter is $q$, and
$p=Q-q=|m|d/3$; indeed $D=kd+md/3$, $h=kd+2md/3$.

When $id\le D$, align $(AC)^{\mathbb Z}$ and $(GC)^{\mathbb Z}$
at the starts of the indicated blocks. The preceding common suffix
is $L=C$. The two right continuations start with $Y$ and then
with $T'$ and $H$, respectively, by \eqref{eq:tail}.
Here $H$ starts $X^\ell$ in the sense that it is a prefix thereof,
and the first $\min(D,kx)$ letters of $T'$ are a prefix of $X^k$.
Since $bd/3=\min(kd,D,h)$, these continuations agree on $F$.

When $id>D$, write $X^i=U^i C$, $X^\ell=B U^\ell$.
Cancelling the outer powers gives $CYYX^kYB=VVU^kV$.
Since $D<id<2D$ (the latter follows from $i+k<2\ell$),
$CY=VZ$, $|Z|=id-D$; since $\ell d>D$, $YB=GV$,
$|G|=\ell d-D$. Cancelling gives $ZYX^kG=VU^k$.
Since $|Z|<D$, $V=ZYJ$, $|J|=2D-id$ and $X^kG=JU^k$.
Now $CY=ZYJZ$ starts with $Z$, and $|C|>|Z|$, so $C=ZT$.
Consequently $TY=YJZ$ and $(GC)Y=Y(BZ)$.
Align $T^{\mathbb Z}$ and $(GC)^{\mathbb Z}$ at the starts of
these blocks. Their common left suffix is $L=T$, of length $D$,
and their right continuations start with $Y$, then $JZ$ and $BZ$.
In this band $i>k$ and $\ell>k$: $2i>k+\ell$ and
$2\ell>i+k$ imply both inequalities by substitution.
Thus $m>0$, $bd/3=kd$. Both $J$ and $B$ have at least this length
since $|J|-kd=(2\ell-i-k)d/3>0$ and $|B|=\ell d$.
The initial $kd$ letters of each are read at the start of a power of $X$.
This completes the construction.

Write $f=|F|$, $n=|L|$, and index the common interval as
$[0,n+y+f)$. We have $f\le q$ in both cases.
The first $r+p$ letters are inside $L$, and each of their first $r$
positions $j$ satisfies $j+Q<n+y+f$. Indeed
\[
 n+f-Q=(a+b-s)d/3\ge (1+1/c)d>x\ge r ,
 \qquad n-p=(n+f-Q)+(q-f)\ge r.
\]
Using period $Q$ and then period $q$ on this common interval gives
the identical letters at $j,j+Q,j+p$. Thus two length-$r$ factors
inside $L\subset X^\omega$ agree at shift $p$. Primitivity implies
$r\mid p$. Since $r\mid x$ and $3p=|m|d$, it follows that
$r\mid |m|u$. But $0<|m|u\le(c+1)u<r$. This contradiction proves
the criterion.
\end{proof}

\begin{theorem}[The entire upper band]\label{tail:band}
Suppose $i+k<2\ell$ and $2i>k+\ell$. Then the tuple forces
periodicity without a length restriction.
\end{theorem}
\begin{proof}
Use the band factorization in the proof of Theorem~\ref{tail:windows},
with $J=V''$. In particular $YB=GZYJ$ and $X^kG=JU^k$.
Put $m=i+\ell-2k>0$, $p=md/3$.
Since $|B|-|J|=p$, suffix comparison gives $B=KJ$, $|K|=p$.
Thus the first $\min(|J|,kx)$ letters of $J$ are read at phases
zero and $p$ of $X^\omega$: one occurrence begins $X^kG$ and
the other is inside $B$, a prefix of $X^\ell$. If $|J|\ge r$,
primitivity gives $r\mid p$ and hence $r\mid mu$. This contradicts
the outer bound $r>(i+\ell-1)u\ge mu>0$ in a counterexample.

It remains to consider $|J|<r\le x$. Write $c=i+\ell-2$.
Both outer exponents exceed $k$, as shown above. Thus $c\ge2$.
We have $|J|=kd+(2\ell-i-k)d/3\ge kd+d/3$.
If $k\ge2$, this exceeds $x$ since $u<d/c$. If $c\ge3$,
even for $k=1$ it exceeds $x$ by the same bound.
The only remaining possibility is $(i,k,\ell)=(2,1,2)$.
Here $d=3e$, $|K|=2e$, $|J|=4e$, and $|J|<x$ gives
$e<u<3e/2$. As $J$ starts $XG$ and is shorter than $X$, write
$X=JP$, $|P|=u-e$. Also $B=KJ$ starts with $J$ since $B$
is a length-$6e$ prefix of $X^2$. Hence $J$ is a length-$4e$
prefix of $KJ$, forcing $J=K^2$.

Let $T$ be the length-$e$ suffix of $K$. Since $X=K^2P$
ends in $U$, we have $U=TP$. Since $X$ also begins with $U$,
$K$ begins with $T$. Its length is $2e$ and it ends with $T$,
so $K=T^2$. Finally $U^2$ is a prefix of $X^2$,
and $2u<3e<4e$, so it is a prefix of $T^4$.
Let $r_0$ denote the primitive root length of $T$. The two copies
of $U$ start at phases $0,u$ in $T^\omega$, and $u>e\ge r_0$.
Thus $r_0\mid u$, and $U$ is a power of the root of $T$.
Since $U=TP$, so is $P$, hence also $X=T^4P$.
This contradicts noncommutation and finishes the band.
\end{proof}

\begin{corollary}[A second low-band comparison]\label{tail:seam}
Suppose $i+k<2\ell$ and $3\le m'=2k-i-\ell\le i+\ell-1$.
Then the tuple forces periodicity.
\end{corollary}
\begin{proof}
Here $i<k$ and $id\le D$ (indeed $i<\ell$ and $i<k$).
Put $p=m'd/3\ge d$. The identities
$X^\ell=HAU^\ell$, $X^kG=T'AU^k$ give two occurrences
of $A U^{\min(k,\ell)}$, starting at $|H|$ and $D$, differing
by $p$. We have $|A|=|G|+p\ge p$ and $kx-D=ku+p\ge u+p$.
Thus at least the first $u+p\ge x$ letters in both occurrences
are inside the indicated powers of $X$. Comparison of a full primitive
rotation gives $r\mid p$. As before $r\mid m'u$ contradicts the
outer bound. Empty images are disposed of by Lemma~\ref{perred}.
\end{proof}

\begin{lemma}[Divisibility obstruction for $2k>i+\ell$]\label{tail:div}
Suppose $i+k<2\ell$, $m'=2k-i-\ell>0$, and $p=m'd/3$.
If $r\mid p$, all four images commute. In particular the preceding seam
corollary holds for every $m'\ge3$, without an upper bound on $m'$.
\end{lemma}
\begin{proof}
Assume $X,U$ do not commute. Here $i<k$, $i<\ell$, $id\le D$,
$D=kd-p<kx$, and $|H|=D-p$. As in the diagonal proof $T'=HK$,
$A=GJ$, $|J|=|K|=p$, and $JCY=CYK$.
The words $K$ and $A$ start at the same phase $|H|$ in $X^\omega$:
use $T'\preceq X^k$ for $K$ and $HA\preceq X^\ell$ for $A$,
where $\preceq$ denotes prefix. Since $r\mid p$ and $J$ is the
length-$p$ suffix of $A$, both $KA$ and $AJ$ read the same periodic
interval starting at $|H|$. Hence $JV=VJ$ by $V=CYA$.
Let $R_1$ be the length-$r$ rotation starting at the end phase of $A$,
namely $\ell d$ modulo $r$; then $J=R_1^{p/r}$ and $V$ is a power
of $R_1$. The factor $U^\ell$ starts at this phase and ends at zero
in $X^\ell$, whence $R_1U^\ell=U^\ell R$. In particular both $J$
and $V$ belong to the group generated by $R,U$.

Also $T'=R^{p/r}H$ because these two prefixes of $X^\omega$ have the
same length and $T'=HK$. From $X^kG=T'AU^k$ and $HA=X^\ell U^{-\ell}$
one obtains
\[
 G=X^{-k}R^{p/r} X^\ell U^{k-\ell}.
\]
Thus $A=GJ$, $C=U^{-i}X^i$ and $Y=C^{-1}VA^{-1}$ are in that
same subgroup. Apply the free-group homomorphism $\nu(R)=0$, $\nu(U)=1$
from Section~\ref{sec:phase-comparisons}. Both $J$ and $V$ have value
zero since $R_1$ is conjugate to $R$ in this subgroup.
Hence $\nu(G)=\nu(A)=k-\ell$, $\nu(C)=-i$ and
$\nu(Y)=i+\ell-k$. The original equation gives
$3(i+\ell-k)=S$, or $i+\ell=2k$, a contradiction.
For $m'\ge3$ the two seam occurrences in the preceding corollary
give $r\mid p$ independently of its former upper bound.
\end{proof}

\begin{corollary}
Every tuple of the requested tail with $k\ge\ell$ forces periodicity.
\end{corollary}
\begin{proof}
Equality is Theorem~\ref{tail:diagonal}. If $k>\ell$, then
$2k-i-\ell=3(k-\ell)+(2\ell-i-k)>3(k-\ell)\ge3$.
Apply Lemma~\ref{tail:div} and its seam consequence.
\end{proof}

\begin{theorem}\label{tail:112full}
The triple \((1,1,2)\) forces periodicity without a length condition.
\end{theorem}
\begin{proof}
Assume \(XU\ne UX\) and write \(d=3e\), \(x=u+3e\). By Lemma~\ref{tail:112} and the corollary following Lemma~\ref{tail:013}, \(u<e\) and \(y<2e-u\). We use \eqref{eq:tail}: \(X=UC\), \(X^2=HAU^2\), \(YXG=ACYAU\), \(YH=GCY\), with \(|A|,|G|,|C|,|H|=e,2e,3e,5e\); \(H=X^\omega[0,5e)\). Write \(c_j\) for the letters of \(X^\omega\), and put \(s=u+y-e\) and \(\delta=2e-y-u\in(0,2e)\).

\emph{Two shifts.} Since \(y<2e=|G|\), the relation \(YH=GCY\) gives \(G=YG_1\), \(H=G_1CY\) with \(|G_1|=2e-y\), so \(C=X^\omega[2e-y,5e-y)\). As also \(C=X^\omega[u,u+3e)\),
\begin{equation}\label{eq:112shift}
c_{u+j}=c_{u+j+\delta}\qquad(0\le j<3e).
\end{equation}
Since \(y<4e\), the relation \(YXG=ACYAU\) makes \(Y\) the prefix of \(AC\) of length \(y\), and \(XG=(AC)[y,4e)\,Y\,AU\). Thus \(X\) is the length-\(x\) prefix of \((AC)[y,4e)\,(AC)^\omega\).

For \(j\ge e+y+u\), the index \(u+j+\delta\) exceeds \(x\); writing \(j'=j+\delta-3e\in[0,\delta)\) and using \(x\)-periodicity, \eqref{eq:112shift} becomes
\begin{equation}\label{eq:112wrap}
c_{j'}=c_{j'+e+y+2u}\qquad(0\le j'<\delta).
\end{equation}

\emph{Case \(s<0\).} Then \(y\le e\) and \(4e-y\ge x\), so \(X=A[y,e)\,C[0,2e+u+y)\). Comparison with \(X=UC\) shows that \(C\) has period \(-s\). For \(j<e+y+u\), \eqref{eq:112shift} stays inside \(C\), and the period \(-s\) gives \(c_{u+j}=c_{u+j+e}\). Fine--Wilf on \(C[0,2e+u+y)\), whose length is at least \(e-s\), gives period \(g=\gcd(e,-s)\), and then all of \(C\) has period \(g\). Here \(g\mid u+y\), so \(e+y+2u\equiv u\pmod g\). By \eqref{eq:112wrap}, \(C[i]=C[i+u]\) on an interval of length \(\delta-u=e-u-s\ge g\), so \(C\) has period \(g'=\gcd(g,u)\). For \(j'<u\), \eqref{eq:112wrap} shows that \(U\) is the factor of \(C\) at offset \(e+y+u\equiv0\pmod{g'}\). Since \(g'\) divides \(u\) and \(3e\), both \(U\) and \(C\) are powers of \(C[0,g')\), so \(XU=UX\).

\emph{Case \(s=0\).} Then \(X=A[y,e)C\), so \(U=A[y,e)\), and \(\delta=e\). By \eqref{eq:112shift}, \(C=\rho^3\) with \(|\rho|=e\). By \eqref{eq:112wrap}, \(X[0,e)=\rho\), so \(\rho=U\rho[0,e-u)\) has period \(u\) and prefix \(U\). The relation \(G=YAU=YX[0,e+u)\) gives \(U=X[e,e+u)\), the suffix of \(\rho\). Hence \(U\) equals its rotation by \(e\bmod u\), and \(U\), \(\rho\) and \(X\) are powers of one word of length \(\gcd(u,e)\).

\emph{Case \(0<s\le u\).} Now \(y\le e\), and \(X=A[y,e)\,C\,Y[0,s)\). Comparison with \(X=UC\) gives period \(s\) on \(X[e-y,x)\). For \(j<2e\), \eqref{eq:112shift} and the period \(s\) give \(c_{u+j}=c_{u+j+\delta+s}=c_{u+j+e}\), since \(\delta+s=e\). So \(C\) has periods \(e\) and \(s\), hence period \(g=\gcd(e,s)\), and so does \(X[e-y,x)\). Let \(P\) be this \(g\)-periodic pattern. Because \(U\) is a suffix of \(X\) (from \(X^2=HAU^2\)) and \(x-u=3e\equiv0\pmod g\), also \(c_j=P(j)\) for \(j<u\). Since \(e-y<u\), every letter of \(X\) is \(P(j)\). As \(g\mid e\) and \(g\mid u+y\), \eqref{eq:112wrap} reads \(P(j')=P(j'+u)\) for \(0\le j'<\delta\). But \(\delta=e-s\) is a positive multiple of \(g\), so this holds for every residue. Hence \(P\), and so \(X\), has period \(\gcd(g,u)\). That number divides \(u\) and \(x\), so \(XU=UX\).

\emph{Case \(s>u\).} Now \(e<y<2e-u\), and \(X=C[y-e,3e)\,Y[0,s)\). Comparison with \(X=UC\) gives period \(s\) on all of \(X\). As in the previous case, \(C\) has periods \(e\) and \(s\), so \(X\) has period \(g=\gcd(e,s)\). Then \eqref{eq:112wrap} with \(\delta=e-s\in g\mathbb{Z}_{>0}\) gives period \(\gcd(g,u)\) as before, and \(XU=UX\).

Each case contradicts \(XU\ne UX\). Lemma~\ref{commouter} completes the proof.
\end{proof}

\subsection{Two local periods}
Throughout this subsection we assume \(i+k<2\ell\), \(2i\le k+\ell\), and \(XU\ne UX\). We use \eqref{eq:tail}, write \(\Pi=X^\omega\) and \(p=|G|-|A|=(i+\ell-2k)d/3\), and recall \(|H|=h=|G|+id\), \(|T'|=D\), \(B=HA\) and \(V=CYA\).

\begin{lemma}\label{tail:twoper}
\begin{enumerate}
\item[(i)] If \(y\le\min(|A|,|G|)\), then \(\Pi[t]=\Pi[t+p]\) for \(|A|-y\le t<kx\), and \(\Pi[t]=\Pi[t+p+y]\) for \(0\le t<|G|-y\).
\item[(ii)] If \(p>0\) and \(|A|\le y<D\), then \(\Pi[t]=\Pi[t+p]\) for \(0\le t<kx\).
\item[(iii)] Let \(\gamma=\gcd(D,h)\). If \(id+y+\min(D,kx,h)\ge D+h-\gamma\), then \(H\) has period \(\gamma\).
\end{enumerate}
\end{lemma}
\begin{proof}
For \(y<\min(D,h)\), the conjugacies \(TY=YT'\) and \(GCY=YH\) give \(AC=Y\tau\), \(T'=\tau Y\), \(GC=Y\tau'\) and \(H=\tau'Y\), where \(\tau=(AC)[y,D)\) and \(\tau'=(GC)[y,h)\). Hence
\[
X^kG=\tau YAU^k\quad(\text{the }k\text{-line}),\qquad X^\ell=\tau'YAU^\ell\quad(\text{the }\ell\text{-line}).
\]
Both \(\tau\) and \(\tau'\) end with a nonempty suffix of \(C\). Let \(\kappa\) be the length of their common suffix. Then \(\kappa YAU^k\) — that is, the common suffix followed by \(YAU^k\) — occurs in the \(k\)-line at position \(D-y-\kappa\) and in the \(\ell\)-line at position \(h-y-\kappa\). These positions differ by \(p\), and the occurrence in the \(\ell\)-line lies inside \(X^\ell\), because it ends at \(\ell d+ku\le\ell x\). The first \(kx\) letters of the \(k\)-line are \(\Pi\).

(i) Here \(\tau=A[y,|A|)C\) and \(\tau'=G[y,|G|)C\), so \(\kappa\ge id\). This gives the first relation. The same alignment applied to the last \(|G|\) letters of the \(k\)-line gives \(G=\Pi[kx+p,kx+p+|G|)=\Pi[p,p+|G|)\). Since \(y\le|G|\), the relation \(YH=GCY\) gives \(Y=G[0,y)\) and \(H=G[y,|G|)\,CY\). Thus \(\Pi[0,|G|-y)=\Pi[p+y,p+|G|)\), which is the second relation.

(ii) Now \(\tau\) is a suffix of \(C\) of length \(D-y\le id\), and \(\tau'\) ends with \(C\) and is longer. So \(\kappa=|\tau|\), and the common factor is the entire \(k\)-line, which occurs at position \(p\) in \(\Pi\).

(iii) Align \((AC)^{\mathbb Z}\) and \((GC)^{\mathbb Z}\) at the starts of the blocks \(AC\) and \(GC\). Both are preceded by \(C\) and followed by \(Y\). After \(Y\), they continue with \(T'\) and with \(H\) respectively. The first \(\min(D,kx)\) letters of \(T'\) and all of \(H\) are prefixes of \(\Pi\). So the two words agree on an interval of length \(id+y+\min(D,kx,h)\). By Fine--Wilf, this interval has period \(\gamma\). It is longer than both periods, so both bi-infinite words have period \(\gamma\), and so does their factor \(H\).
\end{proof}

\begin{lemma}\label{tail:rp}
In a counterexample with \(m=i+\ell-2k\ne0\) and \(|m|\le i+\ell-2\), no integer \(q\) with \(q\mid p\) is a period of a factor of \(\Pi\) of length at least \(x+q\).
\end{lemma}
\begin{proof}
By Fine--Wilf, such a factor has period \(\gcd(q,x)\) and length at least \(x\). Hence \(r\mid q\mid p\), and so \(r\mid 3p=md\). Since \(r\mid mx\), it follows that \(r\mid mu\). But \(0<|m|u<(i+\ell-1)u<r\) by Corollary~\ref{outerbound}, a contradiction.
\end{proof}

\begin{theorem}\label{tail:124}
The triples \((1,2,4)\) and \((2,2,4)\) force periodicity without a length condition.
\end{theorem}
\begin{proof}
Write \(d=3e\). Then \(x=u+3e<4e\), by Corollary~\ref{outerbound} (\(3u<3e\), resp.\ \(4u<3e\)). The table lists \((|A|,|G|,p,D,h,kx)\) for each triple:
\[
(1,2,4):(4e,5e,e,7e,8e,2u+6e),\qquad (2,2,4):(2e,4e,2e,8e,10e,2u+6e).
\]
In both cases \(m\in\{1,2\}\), so Lemma~\ref{tail:rp} applies.

\emph{Case \(y\le|A|\).} By Lemma~\ref{tail:twoper}(i), \(\Pi\) has period \(p\) on \(W_a=[|A|-y,kx+p)\), and period \(p+y\) on \(W_b=[0,|G|+p)=[0,6e)\). Their overlap \([|A|-y,6e)\) has length \(6e-|A|+y=(e,\text{resp. }2e)+p+y\), which is at least \(p+(p+y)\), at least \(p+y\), and at least \(p\). By Fine--Wilf, the overlap has period \(g=\gcd(p,y)\). That period spreads to \(W_a\) and to \(W_b\), since each has a period that \(g\) divides and contains the overlap of length at least that period. It then spreads to their union \([0,2u+6e+p)\). The union has length at least \(x+p\ge x+g\). Lemma~\ref{tail:rp} with \(q=g\) gives a contradiction.

\emph{Case \(|A|<y<D\).} Here \(p>0\), so Lemma~\ref{tail:twoper}(ii) gives period \(p\) on \([0,kx+p)\), which has length at least \(x+p\). Contradiction by Lemma~\ref{tail:rp}.

\emph{Case \(y\ge D\).} We have \(\min(D,kx,h)\ge6e\), so \(id+y+\min(D,kx,h)\ge 3e+7e+6e=16e\ge D+h-\gamma=14e\) in the first case, and \(\ge 6e+8e+6e=20e\ge16e\) in the second. By Lemma~\ref{tail:twoper}(iii), \(H=\Pi[0,h)\) has period \(\gamma=\gcd(D,h)\), which is \(e\) resp.\ \(2e\) and divides \(p\). Since \(h\ge 8e\ge x+\gamma\), Lemma~\ref{tail:rp} gives a contradiction.
\end{proof}

\begin{theorem}\label{tail:general}
Let \(i+k<2\ell\), \(2i\le k+\ell\), \(m=i+\ell-2k\ne0\), \(c=i+\ell-2\ge|m|\), \(c\ge1\). Put
\[
a=k+\ell-2i,\quad g_*=2\ell-i-k,\quad s=i+k+\ell,\quad \eta=2i+2\ell-k,\quad g_0=\gcd(s,\eta),
\]
and assume \((\eta-3-g_0)c\ge3\). Then \((i,0,k,\ell)\) forces periodicity, provided that
\begin{itemize}
\item[(P)] \(m>0\), \(\bigl(3k\ge g_*\) or \(3(k-1)\ge a\bigr)\), and \(\min(s,3k,\eta)\ge g_*-g_0\); or
\item[(N)] \(m<0\), \(3(k-1)\ge a\), and \(\min(s,3k,\eta)\ge a-g_0\).
\end{itemize}
\end{theorem}
\begin{proof}
Put \(f=d/3\). Then \(|A|=af\), \(|G|=g_*f\), \(p=mf\), \(D=sf\), \(h=\eta f\), \(id=3if\) and \(kd=3kf\). Also \(\gamma=\gcd(D,h)=g_0f\) divides \(p=h-D\). In a counterexample, Corollary~\ref{outerbound} gives \(u<3f/c\), so \(h\ge x+\gamma\) by the assumption on \(\eta\). We use Lemmas~\ref{tail:twoper} and~\ref{tail:rp}. Recall \(kx\ge kd\) and \(kx-x\ge(k-1)d\).

\emph{(P).} Here \(p>0\) and \(|A|<|G|\).

\emph{Case \(y\le|A|\).} By Lemma~\ref{tail:twoper}(i), period \(p\) holds on \(W_a=[|A|-y,kx+p)\). If \(3(k-1)\ge a\), then \(|W_a|\ge x+p\), and Lemma~\ref{tail:rp} gives a contradiction. Otherwise \(kd\ge|G|\). Period \(p+y\) holds on \(W_b=[0,|G|+p)\). The overlap \([|A|-y,|G|+p)\) has length \(2p+y\), so it has period \(\gcd(p,y)\), which spreads to \(W_a\cup W_b=[0,kx+p)\), of length at least \(x+p\). Again Lemma~\ref{tail:rp} gives a contradiction.

\emph{Case \(|A|<y<D\).} Lemma~\ref{tail:twoper}(ii) gives period \(p\) on \([0,kx+p)\).

\emph{Case \(y\ge D\).} We have \(id+y+\min(D,kx,h)\ge D+h-\gamma\), because \(\min(D,kd,h)\ge|G|-\gamma\). So \(H\) has period \(\gamma\mid p\), and \(h\ge x+\gamma\).

\emph{(N).} Here \(p<0\) and \(|G|<|A|\).

\emph{Case \(y\le|G|\).} Lemma~\ref{tail:twoper}(i) gives period \(|p|\) on \([|G|-y,kx)\), of length \(kx-|G|+y\ge x+|A|-|G|=x+|p|\).

\emph{Case \(|G|<y<h\).} The proof of Lemma~\ref{tail:twoper}(ii), with the roles of \(\tau\) and \(\tau'\) exchanged, shows that \(\tau'\) is a suffix of \(\tau\). So \(\Pi[t]=\Pi[t+p]\) for \(-p\le t<kx\), that is, period \(|p|\) on \([0,kx)\), and \(kx\ge x+(k-1)d\ge x+|p|\).

\emph{Case \(y\ge h\).} We have \(\min(D,kd,h)\ge|A|-\gamma=D-id-\gamma\), so Lemma~\ref{tail:twoper}(iii) applies again.

In every case Lemma~\ref{tail:rp} gives a contradiction.
\end{proof}

\begin{corollary}\label{tail:gencor}
The hypotheses of Theorem~\ref{tail:general} hold for \((1,2,4)\), \((1,3,4)\), \((1,3,6)\), \((1,4,5)\), \((1,4,6)\), \((2,2,4)\), \((2,2,5)\), \((2,3,6)\), \((3,2,5)\), \((3,2,6)\) and \((3,3,6)\). Hence all eleven force periodicity.
\end{corollary}

\begin{lemma}\label{tail:astar}
For every \(y\), put \(\kappa=\min(D,h,y+id)\). Then \(\Pi[t]=\Pi[t+p]\) for \(D-\kappa\le t<kx\).
\end{lemma}
\begin{proof}
From \(TY=YT'\) and \(GCY=YH\), the word \(T'\) is the suffix of \(ACY\) of length \(D\), and \(H\) is the suffix of \(GCY\) of length \(h\). These two words share the suffix \(CY\), so \(T'\) and \(H\) share a suffix \(\sigma\) of length \(\kappa\). Therefore \(\sigma AU^k\) occurs in \(X^kG=T'AU^k\) at position \(D-\kappa\) and in \(X^\ell=HAU^\ell\) at position \(h-\kappa\), as in the proof of Lemma~\ref{tail:twoper}. The two positions differ by \(p\).
\end{proof}

\begin{theorem}\label{tail:214}
The triples \((2,1,4)\) and \((0,3,4)\) force periodicity without a length condition.
\end{theorem}
\begin{proof}
Write \(d=3e\). In both cases \(|m|=c\), so Lemma~\ref{tail:rp} applies. Note that whenever \(\Pi\) has a global period \(q\mid p\), Lemma~\ref{tail:rp} gives a contradiction.

\emph{\((2,1,4)\).} Here \(|A|=e\), \(|G|=5e\), \(p=4e\), \(D=7e\), \(h=11e\), \(id=6e\), \(kx=x=u+3e\), and \(4u<3e\).

If \(y\ge e\), then \(\kappa=D\), and Lemma~\ref{tail:astar} gives period \(4e\) on \([0,x+4e)\). Fine--Wilf with the period \(x\) gives a global period \(\gcd(4e,x)\mid p\).

If \(y<e\), Lemma~\ref{tail:twoper}(i) gives period \(4e+y\) on \([0,9e)\). Since \(9e\ge(4e+y)+x\), \(\Pi\) has global period \(g_1=\gcd(4e+y,x)\). Because \(x<4e+y<2x\), we have \(g_1\le x/2\le x-e\). Lemma~\ref{tail:astar} gives period \(4e\) on \([e-y,x+4e)\), of length \(x+3e+y\ge4e+g_1\). So \(\Pi\) has global period \(\gcd(4e,g_1)\mid p\).

\emph{\((0,3,4)\).} Here \(C=\varepsilon\), \(|A|=D=7e\), \(|G|=h=5e\), \(p=-2e\), \(kx=3u+9e\), and \(2u<3e\).

By Lemma~\ref{tail:astar} with \(\kappa=\min(5e,y)\), period \(2e\) holds on \([5e-\kappa,kx)\), of length \(3u+4e+\kappa\). If \(2u+\kappa\ge e\), this is at least \(2e+x\), and we obtain a global period \(\gcd(2e,x)\mid p\).

Otherwise \(y<e-2u\). Then Lemma~\ref{tail:twoper}(i) gives period \(2e-y\) on \([y-2e,5e-y)\), of length \(7e-2y\ge(2e-y)+x\). So \(\Pi\) has global period \(g_1=\gcd(2e-y,x)\le2e-y<x\). The first window has length \(3u+4e+y\ge2e+g_1\), so \(\Pi\) has global period \(\gcd(2e,g_1)\mid p\).
\end{proof}

\begin{theorem}\label{tail:023}
The triple \((0,2,3)\) forces periodicity without a length condition.
\end{theorem}
\begin{proof}
Write \(d=3e\). Then \(C=\varepsilon\), \(|A|=D=5e\), \(|G|=h=4e\), \(p=-e\), \(kx=2u+6e\), \(x=u+3e\), \(u<3e\), and \(|m|=c=1\). As before, a global period \(q\mid p\) of \(\Pi\) contradicts Lemma~\ref{tail:rp}.

By Lemma~\ref{tail:astar} with \(\kappa=\min(4e,y)\), period \(e\) holds on \(W_a=[4e-\kappa,kx)\), of length \(2u+2e+\kappa\). If \(u+\kappa\ge2e\), this length is at least \(e+x\), and we obtain the global period \(\gcd(e,x)\mid p\).

Now let \(y<2e-u\), so \(\kappa=y\). Here \(T'=A[y,5e)\,Y\) and \(X^2G=T'AU^2\), so \(\Pi[0,5e-y)=A[y,5e)=\Pi[4e+y,9e)\). Put \(\sigma=4e+y-x=e+y-u\).

If \(\sigma\ne0\), then \(\Pi\) has period \(|\sigma|\) on a window of length \(5e-y+|\sigma|\ge x+|\sigma|\). So \(\Pi\) has global period \(g_1=\gcd(|\sigma|,x)\le|\sigma|\le e+y+u\). As \(|W_a|=2u+2e+y\ge e+g_1\), we obtain the global period \(\gcd(e,g_1)\mid p\).

If \(\sigma=0\), then \(u=e+y\), \(x=4e+y\), and \(y<e/2\). By Lemma~\ref{tail:twoper}(i), period \(e-y\) holds on \(W_b=[y-e,4e-y)\). Its translate \(W_b+x=[3e+2y,8e)\) meets \(W_a=[4e-y,8e+2y)\) in an interval of length at least \(\min(4e+y,5e-2y)\ge2e-y\). By Fine--Wilf the union has period \(\gcd(e,e-y)=\gcd(e,y)\). The union is \([\min(4e-y,3e+2y),\,8e+2y)\), of length \(\min(4e+3y,\,5e)\). This is at least \(x+\gcd(e,y)=4e+y+\gcd(e,y)\): if \(y>0\), then \(\gcd(e,y)\le y<e/2\); if \(y=0\), the length is \(5e=x+e\). So we obtain a global period dividing \(e\).
\end{proof}

\subsection{A certified Fine--Wilf covering of the residual tuples}\label{tail:cert-sec}

This subsection settles the eighteen tuples with exponents at most \(6\) that the preceding results leave open. The proof is computer-assisted: it covers the space of length vectors by finitely many polyhedral cones, and on each cone it exhibits one chain of Fine--Wilf steps that works uniformly. Each cone and its chain is a finite certificate, checked in exact rational arithmetic by a program that re-derives every step independently. We first state the rules, and then the covering.

Throughout, \((i,k,\ell)\) satisfies \(k\ge1\), \(i+k<2\ell\), \(2i\le k+\ell\) and \(m=i+\ell-2k\ne0\), and we consider a counterexample with \(XU\ne UX\), oriented so that \(d>0\). By Lemma~\ref{commouter} it suffices to derive \(XU=UX\) or a contradiction. Put \(\kappa=3\) if \(3\nmid S\) and \(\kappa=1\) otherwise, so that \(d=\kappa w\) for a positive integer \(w\). Every length below is a linear form with integer coefficients in \((w,u,y)\): \(x=\kappa w+u\), \(|C|=id\), \(|A|=D-id\), \(|G|=\ell d-D\), \(|H|=|G|+id\), \(p=md/3\). Let \(\Pi\) be the bi-infinite word with period \(x\) whose factor \(\Pi[0,x)\) is \(X\).

\paragraph{The reduced system.} By \eqref{eq:tail} and the identities preceding it,
\begin{equation}\label{eq:reduced}
X^i=U^iC,\qquad X^\ell=HAU^\ell,\qquad YH=GCY,\qquad YX^kG=ACYAU^k,
\end{equation}
where \(C=\Pi[iu,iu+id)\), \(H=\Pi[0,|H|)\), \(A=\Pi[|H|,\ell d)\), \(U=\Pi[0,u)\) if \(i\ge1\) and \(U=\Pi[-u,0)\) if \(i=0\). The last identity makes \(Y\) the prefix of length \(y\) of \((AC)^\omega\), and the third makes \(G\) the prefix of length \(|G|\) of \(YH\). So both sides of each identity in~\eqref{eq:reduced} are concatenations of factors of \(\Pi\) whose starting positions and lengths are linear forms. Fix a closed or open polyhedral cone \(K\subset\mathbb R^3\) on which the comparisons between the cumulative block lengths of the two sides have constant signs. Aligning the two sides then gives finitely many \emph{relations} \((a,b,L)\), with \(a,b,L\) linear forms: at every integer point of \(K\),
\[
\Pi[a+t]=\Pi[b+t]\qquad(0\le t<L).
\]

\paragraph{Windows and lattices.} For a finite set \(M\) of linear forms, write \(\Lambda(M)\) for the set of integer combinations of \(M\), and \(P_M\) for the greatest common divisor of the values of \(M\) at the point under consideration. A \emph{window} is a triple \((\lambda,\rho,M)\) of linear forms and a finite set of linear forms, meaning that \(\Pi[\lambda,\rho)\) has period \(P_M\). For finite sets \(M_1,M_2\), a \emph{bound} for \(M\) is a form \(\beta(M)\in\Lambda(M)\) that is positive on \(K\), and a \emph{threshold} is
\[
\tau(M_1,M_2)=\begin{cases}
\beta(M_2)&\text{if }\Lambda(M_2)\subseteq\Lambda(M_1),\\
\beta(M_1)&\text{if }\Lambda(M_1)\subseteq\Lambda(M_2),\\
\beta(M_1)+\beta(M_2)-|f|&\text{if }\Lambda(M_1\cup M_2)=\mathbb Z f,\ f\text{ of constant sign on }K,\\
\beta(M_1)+\beta(M_2)&\text{otherwise.}
\end{cases}
\]

\begin{lemma}[Certificate rules]\label{tail:rules}
Let \(K\) be a cone as above, and let all inequalities named below hold on \(K\). At every integer point of \(K\) the following hold.
\begin{enumerate}
\item[(W1)] If \((a,b,L)\) is a relation, \(j\in\mathbb Z\) and \(s=b-a-jx>0\), then \((a,a+L+s,\{s\})\) is a window. If \(s<0\), then \((a+s,a+L,\{-s\})\) is a window.
\item[(W2)] If \((\lambda_1,\rho_1,M_1)\) and \((\lambda_2,\rho_2,M_2)\) are windows, \(j\in\mathbb Z\), \(\lambda_2'=\lambda_2+jx\), \(\rho_2'=\rho_2+jx\), and \(\rho-\lambda'\ge\tau(M_1,M_2)\) for all \(\rho\in\{\rho_1,\rho_2'\}\), \(\lambda'\in\{\lambda_1,\lambda_2'\}\), then \((\lambda,\rho,M_1\cup M_2)\) is a window for all \(\lambda\in\{\lambda_1,\lambda_2'\}\), \(\rho\in\{\rho_1,\rho_2'\}\).
\item[(G)] Suppose \(\Pi\) has period \(P_\Gamma\) for a finite set \(\Gamma\) (initially \(\Gamma=\{x\}\)). If \((\lambda,\rho,M)\) is a window and \(\rho-\lambda\ge\tau(M,\Gamma)\), then \(\Pi\) has period \(P_{\Gamma\cup M}\).
\item[(C)] Let \(Z\) be a set of forms that vanish identically on \(K\). If \(\Pi\) has period \(P_\Gamma\) and \(u\in\Lambda(\Gamma\cup Z)\), then \(XU=UX\). If \(p\in\Lambda(\Gamma\cup Z)\), the point carries no counterexample.
\end{enumerate}
\end{lemma}
\begin{proof}
The value of \(\beta(M)\) is a positive element of the value group generated by \(M\), hence a multiple of \(P_M\), and so \(P_M\le\beta(M)\). Also \(\gcd(P_{M_1},P_{M_2})=P_{M_1\cup M_2}\). In the first case of the threshold \(P_{M_1}\mid P_{M_2}\), so \(P_{M_1}+P_{M_2}-\gcd=P_{M_2}\le\beta(M_2)\); the second case is symmetric. In the third case \(P_{M_1\cup M_2}=|f|\). Hence in all cases \(\tau(M_1,M_2)\ge P_{M_1}+P_{M_2}-\gcd(P_{M_1},P_{M_2})\).

(W1) is the relation read with the period \(x\) of \(\Pi\). For (W2), translation by \(jx\) preserves windows. The hypothesis says that the overlap \(O=[\max\lambda,\min\rho)\) has length at least \(\tau\). By Fine--Wilf, \(\Pi\) restricted to \(O\) has period \(q=P_{M_1\cup M_2}\). Each window has a period that \(q\) divides and contains \(O\), which is at least that period long. So each window is \(q\)-periodic, with the same residue pattern as on \(O\). Hence the union interval \([\min\lambda,\max\rho)\) is \(q\)-periodic, and so is every subinterval named in (W2). (G) is the same argument with \(\Pi\) itself in place of the second window, using \(\rho-\lambda\ge P_\Gamma\). For (C), the forms in \(Z\) have value \(0\), so \(P_\Gamma\) divides the value of every element of \(\Lambda(\Gamma\cup Z)\). If \(P_\Gamma\mid u\), then \(P_\Gamma\mid x\), and \(X\) and \(U\) are powers of the same word of length \(P_\Gamma\), because \(U\) starts at a multiple of \(P_\Gamma\) in \(\Pi\). If \(P_\Gamma\mid p\), Lemma~\ref{tail:rp} applies (all eighteen tuples satisfy \(|m|\le i+\ell-2\)).
\end{proof}

A \emph{certificate} is a cone \(K\) together with a finite sequence of applications of (W1), (W2), (G) ending in (C), all of whose inequalities hold on \(K\). By Lemma~\ref{tail:rules} and homogeneity, it excludes counterexamples at every integer point of \(K\).

\begin{lemma}[Dilation]\label{tail:dilation}
Let \(P_0=(w_0,u_0,y_0)\) be an integer length vector, and suppose that every solution of~\eqref{eq:reduced} at \(P_0\) has \(\Pi\) of period \(g_0\). Then every solution at \(tP_0\), \(t\ge1\), has \(\Pi\) of period \(tg_0\). In particular, if \(g_0\mid u_0\) or \(g_0\mid p_0\), there is no counterexample on the ray through \(P_0\).
\end{lemma}
\begin{proof}
At \(tP_0\) every block of~\eqref{eq:reduced} starts and ends at a multiple of \(t\). So the letter identifications at \(tP_0\) are exactly \(\Pi[ta+r]=\Pi[tb+r]\), \(0\le r<t\), for the identifications \(\Pi[a]=\Pi[b]\) at \(P_0\). The position graph on \(\mathbb Z/tx_0\) is thus the disjoint union over \(r\) of copies of the graph at \(P_0\) under \(v\mapsto tv+r\). Since \(v\) and \(v+g_0\) are joined at \(P_0\), the vertices \(tv+r\) and \(tv+r+tg_0\) are joined at \(tP_0\). The last claim follows as in (C).
\end{proof}

\begin{lemma}[Long \(Y\)]\label{tail:longY}
For each of the eighteen tuples, a counterexample has \(y<Y_*w\) with \(Y_*\) as in Table~\ref{tab:cert}.
\end{lemma}
\begin{proof}
For \((0,1,3)\), Lemma~\ref{tail:013} applies once \(y\ge5e\) (here \(w=e\)), since \(\min(x,4e)\ge3e\). For the other seventeen, put \(\gamma=\gcd(D,|H|)\). Lemma~\ref{tail:twoper}(iii) applies once \(id+y+\min(D,kd,|H|)\ge D+|H|-\gamma\), because \(kx\ge kd\). That is the stated bound. In each case \(\gamma\mid p\) and \(|H|\ge(1+1/c)d+\gamma>x+\gamma\), where \(c=i+\ell-2\), by Corollary~\ref{outerbound}. So Lemma~\ref{tail:rp} gives a contradiction.
\end{proof}

\begin{theorem}\label{tail:certified}
The tuples \((0,1,3)\), \((0,1,4)\), \((1,1,3)\), \((1,1,4)\), \((0,1,5)\), \((0,1,6)\), \((0,2,5)\), \((0,2,6)\), \((0,3,5)\), \((0,4,6)\), \((1,1,5)\), \((1,1,6)\), \((1,2,5)\), \((1,2,6)\), \((2,1,5)\), \((2,1,6)\), \((2,2,6)\) and \((3,1,6)\) force periodicity without a length condition.
\end{theorem}
\begin{proof}
Consider a counterexample with \(XU\ne UX\). By Corollary~\ref{outerbound} and Lemma~\ref{tail:longY}, the point \((u/w,y/w)\) lies in the open polygon
\[
R=\{0<u/w<\kappa/c,\ 0<y/w<Y_*\}.
\]
For each tuple, the file \texttt{checks/cov\(ik\ell\).json} lists finitely many certificates, and finitely many rational points of \(R\) together with the forced period of the reduced system at the primitive integer vector on their ray. The program \texttt{checks/check\_certs.py} verifies the following in exact rational arithmetic, without using the code that produced the data.
\begin{enumerate}
\item For every certificate, it re-derives the relations from~\eqref{eq:reduced} by aligning the blocks itself. It checks that every relation used is one of them (or is a subrelation of one, or holds trivially by \(x\)-periodicity on the cone), and that the sign conditions fixing that alignment hold on the cone. It then checks each application of (W1), (W2), (G) and the final (C). Each stored bound \(\beta\) must lie in the relevant lattice and be positive on the cone, each lattice inclusion or rank-one claim is recomputed, and every inequality is checked on the cone's section \(w=1\) at its vertices. Strict inequalities are handled exactly at vertices and edges.
\item For every listed point, it recomputes the components of the position graph of~\eqref{eq:reduced} at the primitive integer vector by breadth-first search. It checks that the forced period divides \(u_0\) or \(p_0\), and then applies Lemma~\ref{tail:dilation}.
\item It checks that the cones and points cover \(R\). It does this by recursive exact subdivision: a convex piece is split along the defining inequalities of a certificate containing an interior point of the piece, and every piece that is a single point must be one of the listed points.
\end{enumerate}
The counts are in Table~\ref{tab:cert}. Every integer length vector of a counterexample therefore lies in a certified cone or on a certified ray, which is a contradiction.
\end{proof}

\begin{table}[ht]\centering\small
\begin{tabular}{lcccrrrr}
\toprule
\((i,k,\ell)\) & \(\kappa\) & \(u/w<\) & \(Y_*\) & cones & \(u\)/\(p\) closures & on lines & points\\
\midrule
\((0,1,3)\) & 3 & \(3\) & \(5\) & 2283 & 2283/0 & 2 & 0\\
\((0,1,4)\) & 3 & \(3/2\) & \(8\) & 1447 & 1410/37 & 0 & 0\\
\((1,1,3)\) & 3 & \(3/2\) & \(5\) & 984 & 481/503 & 0 & 0\\
\((1,1,4)\) & 1 & \(1/3\) & \(2\) & 1273 & 1273/0 & 0 & 0\\
\((0,1,5)\) & 1 & \(1/3\) & \(3\) & 441 & 441/0 & 0 & 0\\
\((0,1,6)\) & 3 & \(3/4\) & \(14\) & 1409 & 1394/15 & 3 & 0\\
\((0,2,5)\) & 3 & \(1\) & \(8\) & 2308 & 2308/0 & 7 & 0\\
\((0,2,6)\) & 3 & \(3/4\) & \(10\) & 1619 & 391/1228 & 4 & 0\\
\((0,3,5)\) & 3 & \(1\) & \(7\) & 674 & 674/0 & 2 & 0\\
\((0,4,6)\) & 3 & \(3/4\) & \(8\) & 846 & 177/669 & 3 & 0\\
\((1,1,5)\) & 3 & \(3/4\) & \(11\) & 2071 & 1995/76 & 0 & 0\\
\((1,1,6)\) & 3 & \(3/5\) & \(14\) & 1517 & 1511/6 & 0 & 0\\
\((1,2,5)\) & 3 & \(3/4\) & \(7\) & 4448 & 44/4404 & 0 & 0\\
\((1,2,6)\) & 1 & \(1/5\) & \(3\) & 871 & 871/0 & 0 & 0\\
\((2,1,5)\) & 3 & \(3/5\) & \(11\) & 245 & 211/34 & 4 & 1\\
\((2,1,6)\) & 1 & \(1/6\) & \(4\) & 296 & 203/93 & 0 & 0\\
\((2,2,6)\) & 3 & \(1/2\) & \(10\) & 165 & 0/165 & 0 & 0\\
\((3,1,6)\) & 3 & \(3/7\) & \(14\) & 376 & 343/33 & 0 & 0\\
\midrule
total & & & & 23273 & & 25 & 1\\
\bottomrule
\end{tabular}
\caption{Certified coverings for Theorem~\ref{tail:certified}. Here \(d=\kappa w\); the region is \(0<u/w<\kappa/c\), \(0<y/w<Y_*\). Points are closed by Lemma~\ref{tail:dilation}.}\label{tab:cert}
\end{table}


\begin{remark}
The certificates use only Fine--Wilf steps between equal-factor relations of~\eqref{eq:reduced}, so each one is a proof in the style of Theorems~\ref{tail:124}--\ref{tail:023} restricted to its cone. Periods are kept as lattices of linear forms, and the exact Fine--Wilf bound is used whenever the two lattices are nested or their union has rank one. This is what lets the cones close up. Away from these cases, a gcd of incommensurable forms is bounded only by a positive lattice element. The rays \(t(9,4,6)\) and \(t(9,4,7)\) for \((0,1,4)\), on which the simpler window prover of the working files fails, lie inside certified cones. On a lower-dimensional piece a single certificate may use forms that vanish identically on that piece, through the set \(Z\) in (C). This is how the line \(w=u+2y\) for \((0,1,3)\), which accumulates at the corner \(u\to0\), and the line \(y=w+3u\) for \((0,2,5)\) are closed (the column ``on lines'' of Table~\ref{tab:cert}).
\end{remark}

\subsection{Families with a free final exponent}\label{tail:fam-sec}

The covering of Theorem~\ref{tail:certified} treats one tuple at a time. Here we fix \(i\) and \(k\) and let \(\ell\) be arbitrary: one finite set of certificates handles all \(\ell\ge\ell_0\) at once. The point is that for fixed \(i,k\) every length in~\eqref{eq:reduced} is a linear form in \((d,u,y,D)\), because
\[
\ell d=3D-(i+k)d,\quad |C|=id,\quad |A|=D-id,\quad |H|=2D-kd,\quad |G|=2D-(i+k)d,\quad p=D-kd .
\]
The exponent \(\ell\) survives in only two places: the suffix \(U^\ell\) of \(X^\ell\), and the outer bound. We weaken both to linear statements valid for all \(\ell\ge\ell_0\).

\begin{lemma}\label{tail:famy}
Let \(k\ge1\), \(\ell\ge\max(2k-i+1,2i-k)\) and \(i+k<2\ell\). In a counterexample, \(y+id<D\).
\end{lemma}
\begin{proof}
Here \(m=i+\ell-2k>0\), so \(|H|\ge D\). If \(y+id\ge D\), then Lemma~\ref{tail:astar} applies with \(\kappa=D\) and gives \(\Pi[t]=\Pi[t+p]\) for \(0\le t<kx\). This is a factor of length \(kx+p\ge x+p\) with period \(p\), and Lemma~\ref{tail:rp} applies, since \(0<m\le i+\ell-2\).
\end{proof}

\paragraph{The linear system.} Fix \(i\ge0\), \(k\ge1\), \(\ell_0\ge3\) and \(n_0=3\le\ell_0\). In the variables \((d,u,y,D)\), the system \(\Sigma_{i,k}\) consists of the identities \(X^i=U^iC\), \(YH=GCY\) and \(YX^kG=ACYAU^k\) of~\eqref{eq:reduced}, with the blocks given by the forms above. In place of \(X^\ell=HAU^\ell\) it has the statement that \(\Pi[-n_0u,0)=U^{n_0}\). Every relation of \(\Sigma_{i,k}\) holds in every counterexample for the tuple \((i,k,\ell)\) with \(\ell\ge\ell_0\), because \(U^{n_0}\) is a suffix of \(U^\ell\). The parameter region is the cone
\[
R_{i,k,\ell_0}=\{d,u,y>0,\ y+id<D,\ 3D\ge(i+k+\ell_0)d,\ (i+\ell_0-2)u<d\}.
\]
It contains the length vector of every counterexample for \((i,k,\ell)\) with \(\ell\ge\ell_0\), by Corollary~\ref{outerbound} and Lemma~\ref{tail:famy}. In the section \(D=1\) it is a bounded polytope. Its face \(d=0\) corresponds to \(\ell\to\infty\), and the tuple \((i,k,\ell)\) occupies the slice \(d=3/(i+k+\ell)\).

Certificates are as in Lemma~\ref{tail:rules}, now with forms in four variables, and with one more rule.
\begin{enumerate}
\item[(W3)] If \(\Pi\) has period \(P_\Gamma\) and \(j\in\Lambda(\Gamma)\), then translating a window by \(j\) gives a window.
\end{enumerate}
Rule (W3) replaces the translations by multiples of \(x\) in (W2). It lets a certificate align windows whose positions differ by a form such as \(|H|+y\), which is far from any fixed multiple of \(x\) when \(\ell\) is large. Using it, every certificate uses only forms with bounded coefficients. Without it, the covering does not terminate: the cells accumulate at \(d=0\), one slab per multiple of \(x\).

\begin{theorem}\label{tail:families}
For every pair \((i,k)\) in Table~\ref{tab:fam} and every \(\ell\ge\ell_0(i,k)\) as listed there, the tuple \((i,0,k,\ell)\) forces periodicity. In particular, \((0,0,1,\ell)\) forces periodicity for every \(\ell\).
\end{theorem}
\begin{proof}
For each pair, the file \texttt{covIK\_\(i\)\_\(k\).json} (and \texttt{cov01L7.json} for \((0,1)\)), in \texttt{remaining/tail/phase1/checks/}, lists certificates for \(\Sigma_{i,k}\). The program \texttt{check3f.py} verifies them without using the code that produced them. It re-derives the relations of \(\Sigma_{i,k}\), checks every step on the section \(D=1\) of the cone in exact integer arithmetic, and checks that the cones cover \(R_{i,k,\ell_0}\). Any uncovered piece must be a point or a segment. For every admissible slice \(d=3/(i+k+\ell)\) that meets such a piece, the program either shows that the true outer bound \((i+\ell-2)u<d\) excludes the point, or closes it by union-find on the true equation together with Lemma~\ref{tail:dilation}. For the listed pairs, no uncovered pieces occurred.

The second claim combines this with Theorem~\ref{tail:certified} for \(3\le\ell\le6\), Theorem~\ref{rightray} for \(\ell=2\), Theorem~\ref{tail:diagonal} for \(\ell=1\), and the slice \(\ell=0\).
\end{proof}

\begin{table}[ht]\centering\small
\begin{tabular}{lrr}
\toprule
\((i,k)\) & \(\ell_0\) & cones\\
\midrule
\((0,1)\) & 7 & 356\\
\((0,2)\) & 7 & 873\\
\((0,3)\) & 12 & 92\\
\((0,4)\) & 18 & 52\\
\((0,5)\) & 22 & 53\\
\((0,6)\) & 22 & 60\\
\((1,1)\) & 7 & 705\\
\((1,2)\) & 7 & 758\\
\((1,3)\) & 13 & 197\\
\((2,1)\) & 7 & 506\\
\((2,2)\) & 7 & 642\\
\((2,3)\) & 7 & 1011\\
\((2,4)\) & 7 & 850\\
\((3,1)\) & 11 & 236\\
\((3,2)\) & 7 & 659\\
\((4,1)\) & 11 & 311\\
\bottomrule
\end{tabular}
\caption{Families of Theorem~\ref{tail:families}: \((i,0,k,\ell)\) forces periodicity for all \(\ell\ge\ell_0\). Every covering was re-verified by \texttt{check3f.py}.}\label{tab:fam}
\end{table}
